The standard on-off model of temperature control needs to be relatively
inaccurate with a large hysteresis otherwise the life span of circuitry
decreases proportionally with the temperature range.
The ability of the standard on-off model to maintain a set temperature is
severely diminished because of this, and they tend to work better for
realistic thermally isolated systems which don't need a large amount of power
to maintain the set temperature.
However, in this experiment the system is open to its environment and sudden
changes in temperature from external interaction, this is the benefit of PID
control, the system is able to react to these changes and adjust accordingly.
The first step is to make sure the power provided to the heater is
proportional to the amount needed for the set temperature, this helps to
prevent the heater from overshooting the temperature as the power supplied to
the heater decreases as it approaches the set temperature and increases as it
drops below the desired range.
The percentage of power applied is calculated using the error \(E\) as the
difference between the desired and measured temperatures over the width of the
allowed temperature band,

\begin{equation}
    E = \frac{T_{\text{set}} - T_{\text{measured}}}{\text{band}}\label{eq:error}
\end{equation}

which was used to control the amount of power supplied to the heater as a
percentage of the maximum amount of power.

\begin{equation}
    \text{Power} = ( T_{\text{offset}} + k_p E ) P_{\max}\label{eq:power_proportional}
\end{equation}

The factor \(k_p\) was added to control the effects of the error calculation
and control the speed with which the system reacts.
While this modification continues to produce a notable degree of error, it can
be observed that the temperature-time graph follows a decaying sinusoidal
function permits further refinement by integrating over an interval that is
twice the time required for the temperature to increase by an amount
equivalent to the band amplitude.
The integral of the error over this time interval

\begin{equation}
    I = \int {E } dt\label{eq:integral}
\end{equation}

can be used in conjunction with a proportionality constant \(k_i\) to avoid
overcorrecting the power and introducing new unwanted temperature
oscillations.
Adding this factor to equation~\eqref{eq:power_proportional} gives the
following formula for a more steady code for maintaining the set temperature.

\begin{equation}
    \text{Power} = (T_{\text{offset}} + k_p E + k_i I) P_{\max}\label{eq:power_integral}
\end{equation}

The final adjustment is to control the temperature when reacting to external
stimuli using derivative control.
If the temperature of the environment or the block suddenly spikes or drops,
this factor will allow the heater to react proportionally to the rate of
change of the error calculated by equation~\eqref{eq:error}.
This derivative \(D\) can be found by dividing the negative of the change in
error by the time interval over which this change is measured, the action time
of the code, using the following equation.

\begin{equation}
    D = - \frac{dE}{dx} = - \frac{\Delta{E}}{\Delta{x}}\label{eq:derivative}
\end{equation}

This value can be used in conjunction with an associated proportionality
constant \(k_d\) and introduced to equation~\eqref{eq:power_integral} to produce the final
equation for maintaining the set temperature

\begin{equation}
    \text{Power} = ( T_{\text{offset}} + k_p E + k_i I + k_d D) P_{\max}\label{eq:power_PID}
\end{equation}

In this experimental setup it is possible that the derivative control can
cause any noise in the temperature sensor to be amplified so it may be
beneficial to use a low value for the proportionality constant \(k_d\).